\providecommand{\tightlist}{%
  \setlength{\itemsep}{0pt}\setlength{\parskip}{0pt}}
\documentclass{article}
\usepackage{arxiv}
\usepackage[T1]{fontenc}
\usepackage[utf8]{inputenc}
\usepackage{amsmath,amssymb}
\usepackage{graphicx}
\usepackage{hyperref}
\usepackage{libertine}
\usepackage[libertine]{newtxmath}
\usepackage{xcolor}
\usepackage{framed}
\usepackage{booktabs}
\usepackage{CJKutf8}
\definecolor{shadecolor}{RGB}{248,248,248}
\AtBeginDocument{\setlength{\headheight}{20pt}}
\newenvironment{Shaded}{\begin{snugshade}}{\end{snugshade}}

\title{The Irreducible Gap: Why No Discrete Simulation\\of the Yarrow Stalk Oracle Can Reproduce\\the Idealized Probability Distribution}
\author{Augustin Chan\\
\texttt{aug@iterative.day}\\
Independent Researcher}
\date{\today}

\begin{document}
\maketitle

\begin{abstract}
The yarrow stalk method of I-Ching divination produces four line types with idealized probabilities $\frac{1}{16}$, $\frac{5}{16}$, $\frac{7}{16}$, and $\frac{3}{16}$. These probabilities assume that dividing $N$ stalks into two piles yields each mod-4 residue class with exactly probability $\frac{1}{4}$. We show this assumption cannot hold for any finite $N$: dividing 48 stalks into piles of size $1$ through $47$ produces 11 values $\equiv 0 \pmod{4}$ but 12 each for $\equiv 1, 2, 3$, an inherent imbalance that compounds across three rounds to yield a 17.5\% maximum relative error. We further show empirically that constraining the split range to $\pm 32\%$ of the midpoint minimizes this error to $\sim$2.4\%, through a counter-bias cancellation effect across rounds that is not derivable from first principles. Zheng and Cao (2023) observed comparable deviations in a large-scale simulation but did not identify the mechanism. Our analysis provides the first explanation: the gap between idealized and physical yarrow stalk probabilities is an integer boundary effect inherent to dividing finite discrete objects.
\end{abstract}

\textbf{Keywords:} yarrow stalk method, I-Ching, modular arithmetic, discrete probability, simulation

%% ============================================================
\section{Introduction}\label{sec:introduction}
%% ============================================================

The yarrow stalk oracle, described in the \textit{Great Commentary} (\begin{CJK}{UTF8}{gbsn}繫辭傳\end{CJK}) of the \textit{Yijing}, prescribes a ritual in which 50 stalks are used: one is set aside permanently, and the remaining 49 are divided, counted by fours, and recombined across three rounds to produce a single line of a hexagram \cite{wilhelm1950iching}. The standard mathematical analysis derives four line probabilities---$\frac{1}{16}$, $\frac{5}{16}$, $\frac{7}{16}$, $\frac{3}{16}$ for line values 6, 7, 8, 9 respectively---that differ from the symmetric $\frac{1}{8}$, $\frac{3}{8}$, $\frac{3}{8}$, $\frac{1}{8}$ of the three-coin method \cite{shields1999yarrow}. This asymmetry is the mathematical signature of the yarrow stalk method: changing yin (6) is rarer than changing yang (9), and stable yin (8) is more common than stable yang (7).

The derivation rests on a critical assumption: that when $N$ stalks are divided into two piles, each of the four mod-4 residue classes has exactly probability $\frac{1}{4}$. Under this assumption, round~1 (48 stalks) produces $P(\text{large}) = \frac{1}{4}$ and $P(\text{small}) = \frac{3}{4}$, while rounds~2--3 (40 or 44 stalks) produce $P(\text{large}) = P(\text{small}) = \frac{1}{2}$.

We show that this assumption cannot hold for any finite number of stalks, identify the exact mechanism, and characterize the resulting gap.

%% ============================================================
\section{The Yarrow Stalk Process}\label{sec:process}
%% ============================================================

Of the 50 stalks, one is set aside at the outset to represent the Taiji (\begin{CJK}{UTF8}{gbsn}太極\end{CJK}), the unchanging unity underlying all change; it plays no further role. The remaining 49 form the working pile. Each line requires three rounds. In each round:

\begin{enumerate}
\tightlist
\item Remove one stalk between the fingers (the \textit{Wu} stalk).
\item Divide the remaining $N$ stalks into two piles of sizes $L$ and $N - L$, where $1 \leq L \leq N-1$.
\item Count each pile by fours; the remainder $r_i = n_i \bmod 4$, with $0 \mapsto 4$.
\item The round value is $r_L + r_R + 1$ (including the Wu stalk).
\end{enumerate}

Round values can only be 4, 5, 8, or 9. These map to:
\begin{itemize}
\tightlist
\item $4$ or $5 \to$ mapped value 3 (``small'')
\item $8$ or $9 \to$ mapped value 2 (``large'')
\end{itemize}

The sum of three mapped values determines the line: $6 = 2{+}2{+}2$, $7 = 2{+}2{+}3$ (any order), $8 = 2{+}3{+}3$ (any order), $9 = 3{+}3{+}3$.

A ``large'' result occurs when both pile sizes are $\equiv 0 \pmod{4}$, i.e., when the left pile size $L \equiv 0 \pmod{4}$.

%% ============================================================
\section{The Mod-4 Imbalance}\label{sec:imbalance}
%% ============================================================

\subsection{Round 1}

In round~1, 49 stalks minus the Wu stalk leaves $N = 48$ to divide. The left pile $L$ ranges over $\{1, 2, \ldots, 47\}$---a total of 47 values.

The mod-4 distribution across this range:

\begin{table}[h]
\centering
\begin{tabular}{cccc}
\toprule
Residue & Values & Count & Probability \\
\midrule
$\equiv 0$ & $4, 8, 12, \ldots, 44$ & 11 & 23.40\% \\
$\equiv 1$ & $1, 5, 9, \ldots, 45$ & 12 & 25.53\% \\
$\equiv 2$ & $2, 6, 10, \ldots, 46$ & 12 & 25.53\% \\
$\equiv 3$ & $3, 7, 11, \ldots, 47$ & 12 & 25.53\% \\
\bottomrule
\end{tabular}
\caption{Mod-4 residue distribution for $L \in \{1, \ldots, 47\}$.}
\label{tab:mod4-round1}
\end{table}

The idealized model requires each class to have probability exactly $\frac{1}{4} = 25.00\%$. The $\equiv 0$ class falls short by 1.60 percentage points.

\subsection{General statement}

For any integer $N$, the range $\{1, \ldots, N-1\}$ contains $N-1$ values. The mod-4 class counts are:

\begin{equation}
c_k = \left\lfloor \frac{N-1-k}{4} \right\rfloor + \mathbf{1}[k \leq N-1], \quad k \in \{0,1,2,3\}
\end{equation}

These counts are equal (each $\frac{N-1}{4}$) only when $4 \mid (N-1)$, i.e., $N \equiv 1 \pmod{4}$. For $N = 48$, we have $N \equiv 0 \pmod{4}$, so the classes are necessarily unequal.

\subsection{Compounding across rounds}

The round-1 imbalance alone is small ($\sim$1.6~pp). But the stalk counts in rounds~2 and~3 depend on round~1 outcomes (40 or 44 stalks remain), and each has its own mod-4 distribution. The errors compound multiplicatively: three rounds of small biases produce a maximum relative error of 17.5\% in the final line probabilities compared to the idealized distribution.

%% ============================================================
\section{The Constrained Split}\label{sec:constrained}
%% ============================================================

Physical yarrow stalk division is not uniform over $\{1, \ldots, N-1\}$. A person grabbing a bundle would not take just 1 stalk or nearly all of them. We model this with a \textit{variation percentage} parameter $v$ that constrains the left pile to the range:

\begin{equation}
L \in \left[\max\!\left(1,\; \left\lfloor\tfrac{N}{2}\right\rfloor - \left\lfloor\tfrac{Nv}{100}\right\rfloor\right),\;\; \min\!\left(N{-}1,\; \left\lfloor\tfrac{N}{2}\right\rfloor + \left\lfloor\tfrac{Nv}{100}\right\rfloor\right)\right]
\end{equation}

At $v \geq 50$, the full range $\{1, \ldots, N-1\}$ is recovered. We swept $v$ from 10 to 100, generating 500{,}000 lines per value using a seeded Mulberry32 PRNG. Results:

\begin{table}[h]
\centering
\begin{tabular}{rccccr}
\toprule
$v$\% & 6 (Old Yin) & 7 (Young Yang) & 8 (Young Yin) & 9 (Old Yang) & Max Err \\
\midrule
10 & 6.16\% & 35.95\% & 36.78\% & 21.11\% & 15.94\% \\
20 & 5.67\% & 32.66\% & 40.92\% & 20.75\% & 10.68\% \\
28 & 6.55\% & 29.11\% & 44.29\% & 20.06\% & 6.98\% \\
\textbf{32} & \textbf{6.22\%} & \textbf{30.77\%} & \textbf{43.82\%} & \textbf{19.19\%} & \textbf{2.36\%} \\
35 & 6.72\% & 31.70\% & 44.60\% & 16.98\% & 9.43\% \\
40 & 5.82\% & 30.79\% & 42.93\% & 20.45\% & 9.08\% \\
50 & 5.14\% & 28.95\% & 44.71\% & 21.20\% & 17.75\% \\
100 & 5.14\% & 28.95\% & 44.71\% & 21.20\% & 17.75\% \\
\midrule
\textit{Idealized} & \textit{6.25\%} & \textit{31.25\%} & \textit{43.75\%} & \textit{18.75\%} & --- \\
\bottomrule
\end{tabular}
\caption{Simulated line probabilities across variation percentages (500{,}000 lines each).}
\label{tab:sweep}
\end{table}

At $v = 32$, the maximum relative error drops to $\sim$2.4\%, far below the 17.75\% of the full-range split.

\subsection{Why 32\% works}

At $v = 32$, the round-1 split range for 48 stalks is $[9, 39]$---31 values. The $\equiv 0$ class has 7 values (22.6\%), \textit{further} from 25\% than the full range. Yet the overall error is \textit{lower}. This paradox resolves because the constrained range introduces a counter-bias in rounds~2 and~3 (which operate on 40 or 44 stalks with different mod-4 distributions) that partially cancels the round-1 imbalance.

This cancellation is not derivable from first principles. It is an empirical coincidence of how integer boundary effects interact across three rounds with the specific stalk counts prescribed by the ritual (48, then 40 or 44, then further reduced). A different total stalk count would have a different sweet spot.

%% ============================================================
\section{Relation to Prior Work}\label{sec:prior}
%% ============================================================

Zheng and Cao \cite{zheng2023bigdata} conducted a large-scale simulation of the yarrow stalk method, observing line probabilities of 6.387\%, 31.38\%, 43.62\%, and 18.61\%---consistent with our full-range ($v \geq 50$) results and clearly deviating from the idealized distribution. However, they did not identify the mod-4 residue class mechanism as the cause of the deviation, nor did they explore constrained split ranges.

Shields \cite{shields1999yarrow} presented probabilistic studies of the I-Ching using APL but worked within the idealized framework, assuming uniform mod-4 distributions.

The gap we identify is, in a sense, a special case of a general phenomenon in number theory: uniform distributions over $\{1, \ldots, N\}$ do not partition evenly into residue classes unless $N$ is a multiple of the modulus \cite{hardy1979introduction}.

%% ============================================================
\section{Discussion}\label{sec:discussion}
%% ============================================================

The idealized probabilities ($\frac{1}{16}$, $\frac{5}{16}$, $\frac{7}{16}$, $\frac{3}{16}$) describe what the yarrow stalk process would produce if stalks were infinitely divisible or if mod-4 residue classes were perfectly balanced. The simulated probabilities describe what happens when a person divides 49 discrete objects according to the prescribed ritual.

Both are faithful to the yarrow stalk method; they answer different questions. The idealized model answers: \textit{what is the mathematical structure of this divination system?} The simulation answers: \textit{what happens when a person performs it?}

The gap is small in practice ($\sim$2--3\% at the constrained sweet spot, $\sim$17\% at full range) but real, inherent, and---as a property of integer arithmetic---irreducible.

%% ============================================================
\section*{Code Availability}
%% ============================================================

The complete simulation code, analysis scripts, and verification ADR are available at:

\url{https://github.com/augchan42/yarrow-stalk-probabilities}

%% ============================================================
\section*{Acknowledgments}
%% ============================================================

For Dr.\ Bruce Cheung, in memory of Robert Uzgalis and the questions that persist across generations of hardware.

\bibliographystyle{plain}
\bibliography{references}

\end{document}
